% Source: Wald GR, physical PDF pp. 213--223
%         (printed pp. 200--210).
% Coverage: Section 8.3, Domains of Dependence; Global Hyperbolicity;
%           equations (8.3.1)--(8.3.10), theorems 8.3.1, 8.3.5, 8.3.7,
%           8.3.9--8.3.12, 8.3.14, propositions 8.3.2, 8.3.4, 8.3.6,
%           8.3.13, lemmas 8.3.3 and 8.3.8, and unnumbered corollaries.
% Figures/tables/footnotes: Figures 8.10--8.14; footnote 7.
% Status: source-reviewed and compile-clean.
% Uncertainties: none.

% Source: Wald GR, physical PDF p. 213
%         (printed p. 200).
% Coverage: The edge of a closed achronal set.
% Figures/tables/footnotes: Figure 8.10; no tables or footnotes.
% Status: source-reviewed and compile-clean.
% Uncertainties: none.

\begingroup
\renewcommand{\thefigure}{8.10}
\begin{figure}[H]
  \centering
  \begin{tikzpicture}[x=.84cm,y=.84cm,>=Latex,line cap=round,font=\small]
    \draw[thick] (-3.30,.42) .. controls (-2.20,1.12) and (-1.35,.12) .. (-.18,.43)
      .. controls (.52,.61) and (.76,.84) .. (1.14,1.03);
    \draw[thick] (1.17,1.02)
      .. controls (1.35,1.95) and (2.35,2.32) .. (3.05,1.88)
      .. controls (3.78,1.43) and (3.57,.21) .. (2.94,-.22)
      .. controls (2.12,-.76) and (1.27,-.38) .. (1.17,1.02);
    \fill (1.47,1.02) circle (1.55pt) node[left] {\(p\)};
    \fill (1.75,1.69) circle (1.55pt) node[above] {\(q\)};
    \fill (1.55,.28) circle (1.55pt) node[below] {\(r\)};
    \draw[thick] (1.55,.28) .. controls (1.84,.66) and (1.26,1.22) .. (1.75,1.69);
    \draw[->] (2.30,1.35) -- (2.05,1.14);
    \node at (2.47,1.46) {\(\lambda\)};
    \draw[->] (-.85,.78) -- (-.46,.46) node[above=.16cm] {\(S\)};
    \node at (2.96,.92) {\(O\)};
  \end{tikzpicture}
  \caption{A spacetime diagram illustrating the definition of the edge of a closed,
  achronal set \(S\).}
  \label{fig:8.10}
\end{figure}
\endgroup


\subsection{Domains of Dependence; Global Hyperbolicity}\label{sec:8.3}
\setcounter{theorem}{0}

In the previous sections, attention was focused on the collection of events, \(I^+(S)\)
or \(J^+(S)\), that could be influenced by a set, \(S\), of events. In this section, we
shall be concerned with the collection of events which, in the sense described below,
are \lq\lq entirely determined\rq\rq\ by a set of events, \(S\), (which, for technical
reasons, will be taken to be a closed, achronal set). We also shall explore the
properties of spacetimes (or spacetime regions) in which all events are
\lq\lq determined\rq\rq\ by an appropriate \(S\).

We begin, however, by pointing out an important property of closed, achronal sets.
For \(S\) closed and achronal, we define the \emph{edge} of \(S\) as the set of points
\(p\in S\) such that every open neighborhood \(O\) of \(p\) contains a point
\(q\in I^+(p)\), a point \(r\in I^-(p)\) and a timelike curve \(\lambda\) from \(r\)
to \(q\) which does not intersect \(S\) (see Fig.~\ref{fig:8.10}). We often will be
interested in closed, achronal sets without edge. (Such sets are sometimes referred
to as slices.) A repetition of the proof of theorem~\ref{thm:8.1.3} establishes the
following result:

\begin{theorem}\label{thm:8.3.1}
Let \(S\) be a (nonempty) closed achronal set with
\(\operatorname{edge}(S)=\emptyset\). Then \(S\) is a three-dimensional, embedded,
\(C^0\) submanifold of \(M\).
\end{theorem}

Now, let \(S\) be a closed, achronal set (possibly with edge). We define the
\emph{future domain of dependence} of \(S\), denoted \(D^+(S)\), by
\begin{equation}
  D^+(S)=\left\{p\in M\ \middle|\
  \begin{gathered}
    \text{Every past inextendible causal curve}\\
    \text{through }p\text{ intersects }S
  \end{gathered}\right\}.
  \tag{8.3.1}\label{eq:8.3.1}
\end{equation}
Note that we always have \(S\subseteq D^+(S)\subseteq J^+(S)\) and, since \(S\)
is achronal, we also have \(D^+(S)\cap I^-(S)=\emptyset\). Two examples
illustrating the nature of \(D^+(S)\) are given in Figs.~\ref{fig:8.11} and
\ref{fig:8.12}. Our definition of \(D^+(S)\) agrees with that of Hawking and Ellis
(1973) but differs slightly from that of Penrose (1972) and Geroch (1970b) in that
they replace \lq\lq causal curve\rq\rq\ by \lq\lq timelike curve.\rq\rq

The set \(D^+(S)\) is of interest because if \lq\lq nothing can travel faster than light,\rq\rq\
then any signal sent to \(p\in D^+(S)\) must have \lq\lq registered\rq\rq\ on \(S\). Thus,
if we are given appropriate information about \lq\lq initial conditions\rq\rq\ on \(S\), we
should be able to predict what happens at \(p\in D^+(S)\). Conversely, if, say,
\(p\in I^+(S)\) but \(p\notin D^+(S)\), then it should be possible to send a signal
to \(p\) without influencing \(S\), and a knowledge of conditions on \(S\) should not
suffice to determine conditions at \(p\). These expectations are confirmed by an
analysis of the propagation of solutions to hyperbolic wave equations representing
physical fields in curved spacetime (see chapter 10).

% Source: Wald GR, physical PDF p. 214
%         (printed p. 201).
% Coverage: Future domain of dependence and Cauchy horizon with a removed point.
% Figures/tables/footnotes: Figure 8.11; no tables or footnotes.
% Status: source-reviewed and compile-clean.
% Uncertainties: none.

\begingroup
\renewcommand{\thefigure}{8.11}
\begin{figure}[H]
  \centering
  \begin{tikzpicture}[x=.84cm,y=.84cm,>=Latex,line cap=round,font=\small]
    \draw[thick] (-2.80,.34) -- (-1.16,2.95);
    \draw[thick] (-1.16,2.95) -- (2.85,.23);
    \draw[thick] (-2.80,.34) .. controls (-1.50,.70) and (.75,-.20) .. (2.85,.23);
    \node at (-.42,1.45) {\(D^+(S)\)};
    \draw[->] (-1.72,3.18) .. controls (-1.98,2.88) and (-1.99,2.58) .. (-2.05,2.36);
    \draw[->] (.10,3.20) .. controls (.15,2.81) and (.32,2.56) .. (.51,2.30);
    \node at (-.65,3.27) {\(H^+(S)\)};
    \draw[->] (.05,-.11) -- (.08,.26) node[below=.16cm] {\(S\)};
    \draw[fill=white] (.36,2.00) circle (2.2pt);
    \draw[->] (.35,1.76) -- (.35,1.94);
    \node[align=left,right] at (.36,1.59) {(remove\\point)};
  \end{tikzpicture}
  \caption{A spacetime diagram showing the future domain of dependence, \(D^+(S)\),
  and Cauchy horizon \(H^+(S)\) of a particular closed achronal set \(S\) in
  Minkowski spacetime with a point removed.}
  \label{fig:8.11}
\end{figure}
\endgroup


The \emph{past domain of dependence} of \(S\), denoted \(D^-(S)\), is defined by
interchanging \lq\lq future\rq\rq\ and \lq\lq past\rq\rq\ in Eq.~\eqref{eq:8.3.1}. Again, from
a knowledge of conditions on \(S\), we expect to be able to \lq\lq retrodict\rq\rq\
conditions at all \(q\in D^-(S)\). The (full) \emph{domain of dependence} of \(S\),
denoted \(D(S)\), is defined as simply,
\begin{equation}
  D(S)=D^+(S)\cup D^-(S).
  \tag{8.3.2}\label{eq:8.3.2}
\end{equation}
Thus, \(D(S)\) represents the complete set of events for which all conditions should
be determined by a knowledge of conditions on \(S\).

A closed achronal set \(\Sigma\) for which \(D(\Sigma)=M\) is called a
\emph{Cauchy surface}. It follows immediately that for any Cauchy surface \(\Sigma\),
we have \(\operatorname{edge}(\Sigma)=\emptyset\). Thus, by theorem~\ref{thm:8.3.1},
every Cauchy surface is an embedded \(C^0\) submanifold of \(M\). This justifies our
use of the terminology \lq\lq surface\rq\rq\ to describe \(\Sigma\), and, since \(\Sigma\)
is achronal, we may think of \(\Sigma\) as representing an \lq\lq instant of time\rq\rq\
throughout the universe.

A spacetime \((M,g_{ab})\) which possesses a Cauchy surface \(\Sigma\) is said to
be \emph{globally hyperbolic}. (This definition differs significantly from both the
original definition of Leray 1952 and the definition used in Hawking and Ellis 1973,
but all three definitions are equivalent; see the remark at the end of this chapter.)
Thus, in a globally hyperbolic spacetime, the entire future and past history of the
universe can be predicted (or retrodicted) from conditions at the instant of time
represented by \(\Sigma\). Conversely, in a non-globally hyperbolic spacetime we
have a breakdown of predictability in the sense that a complete knowledge of
conditions at a single \lq\lq instant of time\rq\rq\ can never suffice to determine the
entire history of the universe. There are some good reasons for believing that all
physically realistic spacetimes must be globally hyperbolic (see chapter 12 and
Penrose 1979). However, even if one does not accept these arguments and wishes to
consider a non-globally hyperbolic spacetime, one still can apply the theorems proven
below on globally hyperbolic spacetimes to any region of the form
\(\operatorname{int}(D(S))\) for any closed achronal set \(S\).

% Source: Wald GR, physical PDF p. 214
%         (printed p. 201).
% Coverage: A future domain of dependence whose Cauchy horizon intersects \(S\).
% Figures/tables/footnotes: Figure 8.12; no tables or footnotes.
% Status: source-reviewed and compile-clean.
% Uncertainties: none.

\begingroup
\renewcommand{\thefigure}{8.12}
\begin{figure}[H]
  \centering
  \begin{tikzpicture}[x=.83cm,y=.83cm,>=Latex,line cap=round,font=\small]
    \draw[thick] (-3.15,-1.28) -- (-1.06,2.25) -- (2.98,-1.02);
    \draw[thick] (-3.15,-1.28) .. controls (-1.66,1.02) and (-.62,.87) .. (1.17,.15)
      .. controls (1.89,-.17) and (2.40,-.67) .. (2.98,-1.02);
    \node at (-.65,1.16) {\(D^+(S)\)};
    \draw[->] (-1.36,2.58) .. controls (-1.68,2.16) and (-1.75,1.89) .. (-1.78,1.62);
    \draw[->] (.40,2.52) .. controls (.36,2.05) and (.52,1.78) .. (.73,1.48);
    \node at (-.66,2.61) {\(H^+(S)\)};
    \draw[->] (.01,.03) -- (.16,.30) node[below=.10cm] {\(S\)};
  \end{tikzpicture}
  \caption{A spacetime diagram showing \(D^+(S)\) and \(H^+(S)\) for a particular
  closed achronal set \(S\) in Minkowski spacetime. Here \(S\) is \lq\lq asymptotically
  null\rq\rq\ to the right and becomes \lq\lq exactly null\rq\rq\ to the left. This example
  shows that \(H^+(S)\) can intersect \(S\) even if \(\operatorname{edge}(S)=\emptyset\).}
  \label{fig:8.12}
\end{figure}
\endgroup


The closure of the future domain of dependence of \(S\), \(\overline{D^+(S)}\), is
characterized by the following property.

\begin{proposition}\label{prop:8.3.2}
\(p\in\overline{D^+(S)}\) if and only if every past inextendible timelike curve from
\(p\) intersects \(S\).
\end{proposition}
\begin{proof}
If there exists a past inextendible timelike curve from \(p\) which does not intersect
\(S\), it is easy to see that the same property must hold for an open neighborhood
\(O\) of \(p\). In that case, \(O\cap D^+(S)=\emptyset\) and
\(p\notin\overline{D^+(S)}\). Conversely, suppose every past inextendible timelike
curve from \(p\) intersects \(S\). Then either
\(p\in S\subseteq D^+(S)\subseteq\overline{D^+(S)}\) or \(p\in I^+(S)\). If the
latter holds, let \(q\in I^-(p)\cap I^+(S)\) and suppose a past inextendible causal
curve \(\lambda\) from \(q\) failed to intersect \(S\). Then either (i) \(\lambda\)
remains in \(I^+(S)\) or (ii) \(\lambda\) intersects \(\partial I^+(S)\) at a point
\(r\notin S\). In either case, we could construct a past inextendible timelike curve
from \(p\) which does not intersect \(S\). [In case (i), we use lemma~\ref{lem:8.1.4}
to get a timelike curve \(\gamma\subseteq I^+(\lambda)\subseteq I^+(S)\); in case
(ii) we use the corollary to theorem~\ref{thm:8.1.2} to obtain a timelike curve from
\(p\) to \(r\) and then we extend it arbitrarily into the past.] Thus for any
\(q\in I^-(p)\cap I^+(S)\) we have \(q\in D^+(S)\). Since every open neighborhood
of \(p\in I^+(S)\) intersects \(I^-(p)\cap I^+(S)\), we have
\(p\in\overline{D^+(S)}\).
\end{proof}

Some properties of the interiors of \(D^+(S)\) and \(D(S)\) can be seen from the
following lemma, the proof of which is left as an exercise (problem 3).

\begin{lemma}\label{lem:8.3.3}
\begin{align*}
  \operatorname{int}(D^+(S))&=I^-(D^+(S))\cap I^+(S),\notag\\
  \operatorname{int}(D(S))&=I^-(D^+(S))\cap I^+(D^-(S)).
\end{align*}
\end{lemma}

It follows immediately from the definition of a Cauchy surface that if \(\Sigma\) is
a Cauchy surface, then every inextendible causal curve intersects \(\Sigma\). In fact,
we have the following stronger result.

\begin{proposition}\label{prop:8.3.4}
Let \(\Sigma\) be a Cauchy surface and let \(\lambda\) be an inextendible causal
curve. Then \(\lambda\) intersects \(\Sigma\), \(I^+(\Sigma)\), and \(I^-(\Sigma)\).
\end{proposition}
\begin{proof}
Suppose \(\lambda\) did not intersect \(I^-(\Sigma)\). By lemma~\ref{lem:8.1.4} we
could find a past inextendible timelike curve
\(\gamma\subseteq I^+(\lambda)\subseteq I^+(\Sigma\cup I^+(\Sigma))=I^+(\Sigma)\).
If we extend \(\gamma\) indefinitely into the future, it still cannot intersect
\(\Sigma\) or the achronality of \(\Sigma\) would be violated. Since every
inextendible causal curve intersects \(\Sigma\), no such \(\gamma\) exists. Thus
\(\lambda\) must enter \(I^-(\Sigma)\). Similarly, \(\lambda\) must enter
\(I^+(\Sigma)\).
\end{proof}

Let \(S\) be a closed achronal set. We define the \emph{future Cauchy horizon} of
\(S\), denoted \(H^+(S)\), by
\begin{equation}
  H^+(S)=\overline{D^+(S)}-I^-(D^+(S)).
  \tag{8.3.3}\label{eq:8.3.3}
\end{equation}
As will be made more precise later (see the corollary to proposition~\ref{prop:8.3.6}
below), \(H^+(S)\) and the analogously defined \(H^-(S)\) measure the failure of
\(S\) to be a Cauchy surface. Examples of \(H^+(S)\) are shown in
Figs.~\ref{fig:8.11} and \ref{fig:8.12}. Clearly \(H^+(S)\) always is closed since it
is the intersection of the two closed sets \(D^+(S)\) and \(M-I^-(D^+(S))\).
Furthermore, we have
\begin{equation}
  \begin{aligned}
    I^-(H^+(S))&\subseteq I^-(\overline{D^+(S)})\\
      &=I^-(D^+(S))\\
      &\subseteq M-H^+(S).
  \end{aligned}
  \tag{8.3.4}\label{eq:8.3.4}
\end{equation}
Thus, \(I^-(H^+(S))\cap H^+(S)=\emptyset\), so \(H^+(S)\) is achronal. Indeed,
\(H^+(S)\) is a portion of the boundary of the past of the set \(D^+(S)\);
specifically we have
\(H^+(S)=[I^+(S)\cup S]\cap\partial I^-(D^+(S))\) (see problem 5).

One of the most important properties of \(H^+(S)\) is stated in the following
theorem:

\begin{theorem}\label{thm:8.3.5}
Every point \(p\in H^+(S)\) lies on a null geodesic \(\lambda\) contained entirely
within \(H^+(S)\) which either is past inextendible or has a past endpoint on the
edge of \(S\).
\end{theorem}
\begin{proof}
The basic idea of the proof is similar to that of theorem~\ref{thm:8.1.6}. Let
\(p\in H^+(S)\) with \(p\notin\operatorname{edge}(S)\). Then either
(i) \(p\in I^+(S)\) or (ii) \(p\in S\) but \(p\notin\operatorname{edge}(S)\). We
first will show that in either case, a nontrivial past directed null geodesic
contained in \(H^+(S)\) passes through \(p\).

In case (i), since \(p\notin I^-(D^+(S))\), for every \(q\in I^-(p)\) there exists
a past inextendible causal curve from \(q\) which does not intersect \(S\). Let
\(\{q_n\}\) be a sequence of points in \(I^-(p)\) which converges to \(p\), and let
\(\{\lambda_n\}\) be a corresponding sequence of such curves as illustrated in
Fig.~\ref{fig:8.13}. Since \(p\) is a limit point of \(\{\lambda_n\}\), by
lemma~\ref{lem:8.1.5} there exists a past inextendible causal limit curve,
\(\lambda\), of the \(\{\lambda_n\}\) which passes through \(p\). Now, suppose
\(\lambda\) entered the open set \(I^+(S)\cap I^-(D^+(S))\subseteq D^+(S)\).
Then so would some \(\lambda_n\) for sufficiently large \(n\), which is a
contradiction since \(\lambda_n\cap D^+(S)=\emptyset\) because each \(\lambda_n\)
fails to intersect \(S\). In particular, since
\(I^-(p)\subseteq I^-(D^+(S))=I^-(\overline{D^+(S)})\), this implies that within
\(I^+(S)\), \(\lambda\) is a past directed causal curve from \(p\) which does not
enter \(I^-(p)\). Thus, by the corollary to theorem~\ref{thm:8.1.2}, within
\(I^+(S)\), \(\lambda\) must be a null geodesic. Furthermore, if a past directed
timelike curve from a point in \(\lambda\cap I^+(S)\) failed to intersect \(S\), we
could construct a past directed timelike curve from \(p\) with the same property.
Since \(p\in\overline{D^+(S)}\), this is impossible because of
proposition~\ref{prop:8.3.2}, and thus by the same proposition we have
\([\lambda\cap I^+(S)]\subseteq D^+(S)\). Thus,
\(\lambda\cap I^+(S)\subseteq H^+(S)\). Putting all these results together,

% Source: Wald GR, physical PDF p. 217
%         (printed p. 204).
% Coverage: A sequence approaching a Cauchy horizon.
% Figures/tables/footnotes: Figure 8.13; no tables or footnotes.
% Status: source-reviewed and compile-clean.
% Uncertainties: none.

\begingroup
\renewcommand{\thefigure}{8.13}
\begin{figure}[H]
  \centering
  \begin{tikzpicture}[x=.80cm,y=.80cm,>=Latex,line cap=round,font=\small]
    \draw[thick] (-1.55,.12) .. controls (-.55,.58) and (.48,.22) .. (1.25,.34)
      -- (3.05,1.95);
    \draw[thick] (1.25,.34) -- (.13,2.62) -- (3.05,1.95);
    \node at (1.55,1.29) {\(D^+(S)\)};
    \draw[->] (.76,2.90) .. controls (.63,2.52) and (.67,2.26) .. (.84,2.01);
    \draw[->] (1.90,2.91) .. controls (1.96,2.53) and (2.16,2.29) .. (2.39,2.13);
    \node at (1.30,2.97) {\(H^+(S)\)};
    \draw[->] (1.73,.06) -- (1.60,.37) node[below=.16cm] {\(S\)};
    \fill (-.22,1.76) circle (1.6pt) node[right] {\(p\)};
    \foreach \x/\y in {-.42/2.08,-.57/2.37,-.72/2.67,-.87/2.95} {
      \fill (\x,\y) circle (1.2pt);
    }
    \node at (-.61,3.29) {\(q_n\)};
    \draw[thick] (-1.70,-1.10) .. controls (-1.47,.44) and (-.98,1.51) .. (-.42,2.08);
    \draw[thick] (-1.53,-1.00) .. controls (-1.24,.64) and (-.89,1.89) .. (-.57,2.37);
    \draw[thick] (-1.34,-.89) .. controls (-1.05,.78) and (-.83,2.16) .. (-.72,2.67);
    \draw[thick] (-1.16,-.76) .. controls (-.92,.91) and (-.81,2.46) .. (-.87,2.95);
    \draw[->] (-1.08,1.56) -- (-1.59,1.31);
    \node at (-1.68,1.79) {\(\lambda_n\)};
  \end{tikzpicture}
  \caption{A point \(p\in H^+(S)\cap I^+(S)\) with a sequence of points
  \(q_n\in I^+(p)\) converging to \(p\) (see theorem~\ref{thm:8.3.5}).}
  \label{fig:8.13}
\end{figure}
\endgroup


starting from every \(p\in H^+(S)\cap I^+(S)\) we have obtained a nontrivial past
directed null geodesic segment \(\lambda\) lying in \(H^+(S)\).

In case (ii) where \(p\in S\) but \(p\notin\operatorname{edge}(S)\), we use the
definition of \(\operatorname{edge}(S)\) to establish existence of an open set \(O\)
such that no causal curve contained in \(O\) from a point
\(q\in I^-(p)\cap O\) can enter \(I^-(S)\cap O\) without intersecting \(S\). The
argument establishing existence of a nontrivial past directed null geodesic through
\(p\) remaining in \(H^+(S)\) then proceeds along similar lines.

Finally, suppose a past directed null geodesic \(\lambda\) leaves \(H^+(S)\); i.e.,
suppose the portion of \(\lambda\) contained in \(H^+(S)\) has a past endpoint \(r\).
Since \(H^+(S)\) is closed, we have \(r\in H^+(S)\). If
\(r\notin\operatorname{edge}(S)\), we can find a nontrivial past directed null
geodesic segment \(\lambda'\) in \(H^+(S)\) starting from \(r\). However, if
\(\lambda'\) were not a continuation of \(\lambda\), by the corollary to
theorem~\ref{thm:8.1.2} we could find a timelike curve connecting a point of
\(\lambda\) to a point of \(\lambda'\). This would violate the achronality of
\(H^+(S)\). This proves that \(\lambda\) cannot have a past endpoint except on
\(\operatorname{edge}(S)\).
\end{proof}

The (full) \emph{Cauchy horizon} of a closed achronal set \(S\) is defined by
\begin{equation}
  H(S)=H^+(S)\cup H^-(S),
  \tag{8.3.5}\label{eq:8.3.5}
\end{equation}
where \(H^-(S)\) is defined by interchanging past and future in the definition of
\(H^+(S)\). Using lemma~\ref{lem:8.3.3}, it is not difficult to show that the Cauchy
horizon marks the boundary of the domain of dependence of \(S\):

\begin{proposition}\label{prop:8.3.6}
\(H(S)=\partial D(S)\).
\end{proposition}

We leave the proof of this proposition as an exercise (problem 6). As a direct
corollary we have

\medskip
\noindent\textsc{Corollary.}\quad If \(M\) is connected, then a nonempty closed
achronal set, \(\Sigma\), is a Cauchy surface for \((M,g_{ab})\) if and only if
\(H(\Sigma)=\emptyset\).
\medskip
\begin{proof}
If \(\partial D(\Sigma)=\emptyset\), we have
\(\overline{D(\Sigma)}=\operatorname{int}(D(\Sigma))=D(\Sigma)\), so \(D(\Sigma)\)
is both open and closed. Thus, since \(D(\Sigma)\supseteq\Sigma\ne\emptyset\) and
\(M\) is connected, we have \(D(\Sigma)=M\).
\end{proof}

This corollary leads to the following useful criterion for Cauchy surfaces.

\begin{theorem}\label{thm:8.3.7}
If \(\Sigma\) is a closed, achronal, edgeless\footnote{The requirement that
\(\Sigma\) be edgeless can be removed from the hypothesis of this theorem. See
Geroch (1970b).} set, then \(\Sigma\) is a Cauchy surface if and only if every
inextendible null geodesic intersects \(\Sigma\) and enters \(I^+(\Sigma)\) and
\(I^-(\Sigma)\).
\end{theorem}
\begin{proof}
The \lq\lq only if\rq\rq\ part is a special case of proposition~\ref{prop:8.3.4}. To
prove the \lq\lq if\rq\rq\ part, it suffices to show that if \(\Sigma\) is not a Cauchy
surface, then at least one null geodesic fails to enter \(I^+(\Sigma)\) or
\(I^-(\Sigma)\). But if \(\Sigma\) fails to be a Cauchy surface, then at least one of
\(H^+(\Sigma)\) or \(H^-(\Sigma)\) is nonempty (unless \(M\) is disconnected and
\(\Sigma\) does not intersect a component of \(M\), in which case the result is
trivial). If, say, \(H^+(\Sigma)\ne\emptyset\), then since
\(\operatorname{edge}(\Sigma)=\emptyset\), by theorem~\ref{thm:8.3.5} there exists
a past inextendible null geodesic which remains forever in \(H^+(\Sigma)\) and thus
never enters \(I^-(\Sigma)\). Clearly, if we extend this geodesic forever into the
future, it still cannot enter \(I^+(\Sigma)\) or the achronality of \(\Sigma\) would
be violated.
\end{proof}

Now, let \((M,g_{ab})\) be a globally hyperbolic spacetime with Cauchy surface
\(\Sigma\). It is easy to see that no closed timelike curves can exist in \(M\): A
closed timelike curve which intersects \(\Sigma\) would violate achronality of
\(\Sigma\), whereas a closed timelike (or even causal) curve which fails to intersect
\(\Sigma\) would violate global hyperbolicity, since we could follow this curve
\lq\lq around and around\rq\rq\ to define an inextendible causal curve which does not
intersect \(\Sigma\). In fact, we will see later (theorem~\ref{thm:8.3.14} below)
that \((M,g_{ab})\) must be stably causal. However, we first show that strong
causality holds:

\begin{lemma}\label{lem:8.3.8}
Let \((M,g_{ab})\) be a globally hyperbolic spacetime. Then \((M,g_{ab})\) is
strongly causal.
\end{lemma}
\begin{proof}
In a globally hyperbolic spacetime with Cauchy surface \(\Sigma\), we clearly have
\(M=I^+(\Sigma)\cup\Sigma\cup I^-(\Sigma)\). Suppose strong causality were violated
at \(p\in I^+(\Sigma)\). It follows directly from the definition of strong causality
violation that we could find a convex normal neighborhood \(U\) of \(p\) contained in
\(I^+(\Sigma)\) and a nested family of open sets \(O_n\subseteq U\) which converges
to \(p\) such that for each \(n\) we can find a future directed timelike curve
\(\lambda_n\) which begins in \(O_n\), leaves \(U\), and ends in \(O_n\). Using
lemma~\ref{lem:8.1.5}, we can find a limit curve \(\lambda\) which passes through
\(p\). Although each of the \(\lambda_n\) are extendible, it is clear that
\(\lambda\) must either be inextendible or yield a closed causal curve through \(p\),
in which case it could be made to be inextendible by going \lq\lq around and around.\rq\rq\
Since none of the \(\lambda_n\) can enter \(I^-(\Sigma)\) or achronality of
\(\Sigma\) would be violated, \(\lambda\) also cannot enter \(I^-(\Sigma)\).
However, this contradicts proposition~\ref{prop:8.3.4}. Thus, strong causality
cannot be violated at \(p\in I^+(\Sigma)\) or, by the same reasoning, at
\(p\in I^-(\Sigma)\). For the case \(p\in\Sigma\), we can choose the \(\{O_n\}\)
so that any future directed timelike curve starting in \(O_n\) must exit \(O_n\) in
\(I^+(\Sigma)\). Again we would find that the limit curve \(\lambda\) could not
enter \(I^-(\Sigma)\), in contradiction to proposition~\ref{prop:8.3.4}.
\end{proof}

The remaining results of this section deal with the space of causal curves joining
two points in a globally hyperbolic spacetime. The arguments used in the proofs
provide a beautiful illustration of the abstract topological space methods outlined in
appendix A.

Let \((M,g_{ab})\) be a strongly causal spacetime and let \(p,q\in M\). We define
\(C(p,q)\) to be the set of continuous, future directed causal curves from \(p\) to
\(q\), where curves that differ only by reparameterization are considered to be the
same curve. [Of course, \(C(p,q)\) will be empty unless \(q\in J^+(p)\).] We define a
topology, \(\mathcal T\), on \(C(p,q)\), thereby making \((C(p,q),\mathcal T)\) into a
topological space as follows. Let \(U\subseteq M\) be open, and define
\(O(U)\subseteq C(p,q)\) by
\begin{equation}
  O(U)=\{\lambda\in C(p,q)\mid\lambda\subseteq U\}.
  \tag{8.3.6}\label{eq:8.3.6}
\end{equation}
In other words, \(O(U)\) consists of all causal curves from \(p\) to \(q\) which lie
entirely within \(U\). We define our topology \(\mathcal T\) by calling a subset,
\(\mathcal O\), of \(C(p,q)\) open if it can be expressed as
\begin{equation}
  \mathcal O=\bigcup O(U),
  \tag{8.3.7}\label{eq:8.3.7}
\end{equation}
where each \(O(U)\) is of the form Eq.~\eqref{eq:8.3.6}.

Since \(O(U_1)\cap O(U_2)=O(U_1\cap U_2)\), it is easy to check that \(\mathcal T\)
does indeed define a topology on \(C(p,q)\). Since we restrict consideration to
spacetimes in which no closed causal curves exist, it follows that if
\(\lambda,\lambda'\in C(p,q)\) are distinct causal curves, the subset \(\lambda\) of
\(M\) cannot contain or be contained in the subset \(\lambda'\) of \(M\). From the
properties of the topology on \(M\), it then follows that the topology \(\mathcal T\)
on \(C(p,q)\) is Hausdorff. Furthermore, when no closed causal curves exist,
\(\mathcal T\) is second countable. (See Geroch 1970b for a sketch of the proof of
this result.) Finally, the notion of convergence defined by \(\mathcal T\) is the
following: \(\lambda_n\to\lambda\) if for every open set \(U\subseteq M\) with
\(\lambda\subseteq U\), there exists an \(N\) such that \(\lambda_n\subseteq U\) for
all \(n>N\). If strong causality holds in \(M\) as we require, one can show that this
notion of convergence is equivalent to the notion of convergence of curves defined
in section 8.1. Similarly, the notion of a \lq\lq limit curve,\rq\rq\ \(\lambda\), of a
sequence \(\{\lambda_n\}\) in the sense of section 8.1 coincides with the topological
space notion of an accumulation point defined in appendix A.

The key theorem upon which all further results of this section are based is the
following:

\begin{theorem}\label{thm:8.3.9}
Let \((M,g_{ab})\) be a globally hyperbolic spacetime and let \(p,q\in M\). Then
\(C(p,q)\) is compact.
\end{theorem}
\begin{proof}
Since the topology on \(C(p,q)\) is second countable, by theorem A.9 of appendix A
we need only show that every infinite sequence \(\{\lambda_n\}\) of points (i.e.,
curves) in \(C(p,q)\) has an accumulation point (i.e., a limit curve, \(\lambda\)) in
\(C(p,q)\). Consider, first, the case \(p,q\in D^-(\Sigma)\), where \(\Sigma\) is a
Cauchy surface for \((M,g_{ab})\), and let \(\{\lambda_n\}\) be a sequence in
\(C(p,q)\) as illustrated in Fig.~\ref{fig:8.14}. If we (temporarily) remove \(q\)
from \(M\), then \(\{\lambda_n\}\) becomes a sequence of future inextendible causal
curves starting at \(p\). Hence, by lemma~\ref{lem:8.1.5}, there exists a future
inextendible (in \(M-q\)) limit curve \(\lambda\) starting at \(p\). Since none of the
\(\lambda_n\) enters \(I^-(\Sigma)\), neither can \(\lambda\). If, now, we restore the
point \(q\), then in \(M\) either (i) \(\lambda\) will remain inextendible or (ii)
\(q\) will be an endpoint of \(\lambda\). However, possibility (i) is ruled out by
proposition~\ref{prop:8.3.4} since \(\lambda\) does not enter \(I^-(\Sigma)\).
Hence, \(\lambda\) (with its endpoint at \(q\) added) provides the desired limit
curve.

% Source: Wald GR, physical PDF p. 220
%         (printed p. 207).
% Coverage: A sequence of causal curves from \(p\) to \(q\).
% Figures/tables/footnotes: Figure 8.14; no tables or footnotes.
% Status: source-reviewed and compile-clean.
% Uncertainties: none.

\begingroup
\renewcommand{\thefigure}{8.14}
\begin{figure}[H]
  \centering
  \begin{tikzpicture}[x=.86cm,y=.86cm,>=Latex,line cap=round,font=\small]
    \draw[thick] (-2.65,3.08) .. controls (-1.49,3.83) and (-.24,2.82) .. (1.10,3.18)
      .. controls (1.87,3.36) and (2.36,3.31) .. (3.06,3.48);
    \node at (1.45,3.47) {\(\Sigma\)};
    \fill (0,.06) circle (1.65pt) node[below] {\(p\)};
    \fill (0,2.26) circle (1.65pt) node[above] {\(q\)};
    \draw[thick] (0,.06) .. controls (-.64,.88) and (-.58,1.75) .. (0,2.26);
    \draw[thick] (0,.06) .. controls (-.23,.87) and (.40,1.37) .. (0,2.26);
    \draw[thick] (0,.06) .. controls (.57,.89) and (.62,1.73) .. (0,2.26);
    \draw[->] (.77,1.54) -- (.20,1.50);
    \node at (1.13,1.67) {\(\lambda_n\)};
  \end{tikzpicture}
  \caption{A sequence \(\{\lambda_n\}\) of causal curves from \(p\) to \(q\) in
  the case \(p,q\in D^-(\Sigma)\) (see theorem~\ref{thm:8.3.9}).}
  \label{fig:8.14}
\end{figure}
\endgroup


The case \(p,q\in D^+(\Sigma)\) obviously follows by the same argument. Thus, the
only remaining nontrivial case is \(p\in D^-(\Sigma)\), \(q\in I^+(\Sigma)\). Given
a sequence \(\{\lambda_n\}\) in \(C(p,q)\), the above argument proves existence of a
future directed limit curve \(\lambda\) starting at \(p\) which enters \(I^+(\Sigma)\).
We choose \(r\in\lambda\cap I^+(\Sigma)\) and extract a subsequence
\(\{\lambda'_n\}\) such that every point on the segment of \(\lambda\) between \(p\)
and \(r\) is a convergence point of this subsequence. Now we reverse the procedure
and consider the sequence of past inextendible (in \(M-p\)) causal curves
\(\{\lambda'_n\}\) starting from \(q\). By the same arguments as given above, we
obtain a limit curve \(\lambda'\) from \(q\) which enters \(I^-(\Sigma)\) and thus
must pass through \(r\), since \(r\) is a convergence point of \(\{\lambda'_n\}\) and
if \(\lambda'\) did not extend to \(r\) it would have to remain in
\(I^+(r)\subseteq I^+(\Sigma)\). Thus, by joining the segment of \(\lambda'\) from
\(r\) to \(q\) with the segment of \(\lambda\) from \(p\) to \(r\) we get the desired
limit curve.
\end{proof}

The compactness of \(C(p,q)\) directly yields a corresponding compactness of the
subset \(J^+(p)\cap J^-(q)\) in the manifold topology:

\begin{theorem}\label{thm:8.3.10}
Let \((M,g_{ab})\) be a globally hyperbolic spacetime and let \(p,q\in M\). Then
\(J^+(p)\cap J^-(q)\) is compact.
\end{theorem}
\begin{proof}
Since paracompactness implies that the manifold topology is second countable (see
appendix A), by theorem A.9 in that appendix we need only show that every sequence
\(\{r_n\}\) of points in \(J^+(p)\cap J^-(q)\) has an accumulation point \(r\). To
prove this, we let \(\{\lambda_n\}\) be a sequence of causal curves from \(p\) to
\(q\) such that \(\lambda_n\) passes through \(r_n\). Since \(C(p,q)\) is compact and
first countable, by theorem A.9 we find a subsequence \(\{\lambda'_n\}\) which
converges to a curve \(\lambda\in C(p,q)\). Now view \(\lambda\) as a subset of
\(M\). Since \(\lambda\) is compact (for it is the continuous image of a closed
interval of \(\mathbb{R}\)), we can find an open neighborhood \(U\subseteq M\) of
\(\lambda\) such that \(\overline U\) is compact. (Proof. Cover \(\lambda\) with
open sets with compact closure, use compactness of \(\lambda\) to extract a finite
subcover, and take the union.) By definition of convergence, there exists an integer
\(N\) such that for all \(n>N\) we have \(\lambda'_n\subseteq U\). Thus, since
\(r'_n\in\lambda'_n\) and \(U\subseteq\overline U\), we have a subsequence
\(\{r'_n\}\) contained in \(\overline U\). From the compactness of \(\overline U\),
there exists an accumulation point \(r\in\overline U\). If \(r\notin\lambda\), we
would contradict the fact that \(\lambda\) is the limit of \(\{\lambda'_n\}\). Thus,
we have \(r\in\lambda\subseteq J^+(p)\cap J^-(q)\) and \(r\) is the desired
accumulation point of \(\{r_n\}\).
\end{proof}

From theorem A.2 in appendix A, we obtain the following corollary:

\medskip
\noindent\textsc{Corollary.}\quad In a globally hyperbolic spacetime,
\(J^+(p)\cap J^-(q)\) is closed.
\medskip

In fact, it is not difficult to see that \(J^+(p)\) itself must be closed. Namely, if
not we could find a point \(r\in\overline{J^+(p)}\) with \(r\notin J^+(p)\). Choose
\(q\in I^+(r)\). Then we would have \(r\in\overline{J^+(p)}\cap J^-(q)\) but
\(r\notin J^+(p)\cap J^-(q)\), which is a contradiction since
\(J^+(p)\cap J^-(q)\) is closed. One can strengthen these arguments to obtain the
following more general result, the proof of which is left as an exercise (problem 7).

\begin{theorem}\label{thm:8.3.11}
Let \((M,g_{ab})\) be a globally hyperbolic spacetime and let \(K\subseteq M\) be
compact. Then \(J^+(K)\) is closed.
\end{theorem}

Recall that for any subset \(S\), we have
\(I^+(S)\subseteq J^+(S)\subseteq\overline{I^+(S)}\). Thus,
theorem~\ref{thm:8.3.11} shows that for a compact set \(K\) in a globally
hyperbolic spacetime, we have \(J^+(K)=\overline{I^+(K)}\). This implies that
\(\partial I^+(K)\subseteq J^+(K)\), and thus in this case we can strengthen
theorem~\ref{thm:8.1.6} to conclude that every \(p\in\partial I^+(K)\) can be
joined to \(K\) by a past directed null geodesic lying in \(\partial I^+(K)\). Thus,
the phenomenon illustrated in Fig.~\ref{fig:8.2} cannot occur for compact sets in a
globally hyperbolic spacetime.

Let \((M,g_{ab})\) be a globally hyperbolic spacetime with Cauchy surface
\(\Sigma\). For \(q\in D^+(\Sigma)\) we define \(C(\Sigma,q)\) as the set of
continuous future directed causal curves from \(\Sigma\) to \(q\) and we define a
topology on \(C(\Sigma,q)\) in the same way as for \(C(p,q)\). Essentially the same
argument as given in theorem~\ref{thm:8.3.9} establishes that \(C(\Sigma,q)\) also
is compact. A repetition of the proof of theorem~\ref{thm:8.3.10} then yields the
following theorem.

\begin{theorem}\label{thm:8.3.12}
Let \((M,g_{ab})\) be a globally hyperbolic spacetime with Cauchy surface
\(\Sigma\) and let \(q\in D^+(\Sigma)\). Then \(J^+(\Sigma)\cap J^-(q)\) is compact.
\end{theorem}

Our final results concern the topology and causal structure of globally hyperbolic
spacetimes.

\begin{proposition}\label{prop:8.3.13}
Suppose \(\Sigma\) and \(\Sigma'\) both are Cauchy surfaces for the spacetime
\((M,g_{ab})\). Then \(\Sigma\) and \(\Sigma'\) are homeomorphic.
\end{proposition}
\begin{proof}
By lemma~\ref{lem:8.1.1}, there exists a nowhere vanishing timelike vector field
\(t^a\) on \(M\). Since \(\Sigma\) and \(\Sigma'\) are Cauchy surfaces, each
integral curve of \(t^a\) must have precisely one intersection point with \(\Sigma\)
and \(\Sigma'\). The map of \(\Sigma\) onto \(\Sigma'\) obtained by associating
points which lie on the same integral curve of \(t^a\) yields a continuous,
one-to-one map with continuous inverse.
\end{proof}

Our final theorem greatly strengthens lemma~\ref{lem:8.3.8}.

\begin{theorem}\label{thm:8.3.14}
Let \((M,g_{ab})\) be a globally hyperbolic spacetime. Then \((M,g_{ab})\) is
stably causal. Furthermore, a global time function, \(f\), can be chosen such that
each surface of constant \(f\) is a Cauchy surface. Thus \(M\) can be foliated by
Cauchy surfaces and the topology of \(M\) is \(\mathbb{R}\times\Sigma\), where
\(\Sigma\) denotes any Cauchy surface.
\end{theorem}
\begin{proof}[Sketch of proof]
As in theorem~\ref{thm:8.2.2}, we introduce a volume measure \(\mu\) on \(M\) such
that \(\mu(M)\) is finite, and we define the function \(f^-\) by
\begin{equation}
  f^-(p)=\mu(I^-(p)).
  \tag{8.3.8}\label{eq:8.3.8}
\end{equation}
Using the fact that, for all \(q\in M\), \(J^+(q)\) and \(J^-(q)\) are closed, it
is possible to show that \(f^-(p)\) is continuous (i.e., the \lq\lq averaging\rq\rq\
used in theorem~\ref{thm:8.2.2} is not needed); see Hawking and Ellis (1973) or
Geroch (1970b) for details. Thus, \(f^-\) defines a continuous global time function,
which can be \lq\lq smoothed\rq\rq\ to yield a differentiable global time function.
This proves stable causality. To obtain the stronger conclusions of this theorem, we
define \(f^+\) by
\begin{equation}
  f^+(p)=\mu(I^+(p)),
  \tag{8.3.9}\label{eq:8.3.9}
\end{equation}
and \(f\) by
\begin{equation}
  f(p)=f^-(p)/f^+(p).
  \tag{8.3.10}\label{eq:8.3.10}
\end{equation}
Then \(f\) is continuous, and each surface of constant \(f\) is a closed achronal
set. It remains to show that every inextendible causal curve intersects each surface
of constant \(f\). This must occur if \(f^-(p)\) goes to zero along every past
inextendible causal curve, and if, similarly, \(f^+(p)\) goes to zero along every
future inextendible causal curve, since then \(f\) will attain all values in the
interval \((0,\infty)\) along every inextendible causal curve. To prove this, we
suppose that \(\lambda\) is a past inextendible causal curve starting at \(q\) along
which \(f^-\) does not approach zero in the past. Then there must exist a point
\(p\) such that \(p\in I^-(r)\) for all \(r\in\lambda\). Hence \(\lambda\) must be
contained within the set \(J^-(q)\cap J^+(p)\). However, by
theorem~\ref{thm:8.3.10} this set is compact, and by lemma~\ref{lem:8.3.8} strong
causality holds. Thus, by lemma~\ref{lem:8.2.1}, \(\lambda\) must have a past
endpoint, which contradicts the fact that \(\lambda\) is past inextendible. Thus,
\(f^-\) must approach zero as one goes along \(\lambda\) into the past. Clearly,
the analogous property also holds for \(f^+\).
\end{proof}

\medskip
\noindent\emph{Remark.} Our definition of global hyperbolicity of \((M,g_{ab})\)
is (I) existence of a Cauchy surface. Leray's (1952) original definition was
essentially (II) strong causality holds and \(C(p,q)\) is compact for all
\(p,q\in M\). The definition given by Hawking and Ellis (1973) is (III) strong
causality holds and \(J^+(p)\cap J^-(q)\) is compact. Lemma~\ref{lem:8.3.8} and
theorem~\ref{thm:8.3.9} prove that (I) \(\Rightarrow\) (II), while
theorem~\ref{thm:8.3.10} shows (II) \(\Rightarrow\) (III) since the proof used
only the compactness of \(C(p,q)\). Finally, the proof of theorem~\ref{thm:8.3.14}
establishes that (III) \(\Rightarrow\) (I), since only properties derived from
(III) played a role in the construction of Cauchy surfaces. Thus, all three
definitions are equivalent.
